
Documentation frameworks for ML datasets -- including Datasheets for Datasets \citep{Gebru2021Datasheets}, Data Statements \citep{Bender2018DataStatements}, and Data Cards \citep{Pushkarna2022DataCards} -- specify what information dataset creators should record. The Data Transparency Scorecard (DTS; Rondina et al., 2025) operationalizes this into a six-section rubric with inverse-frequency weights, where sections that are rarely documented (e.g., Preprocessing, Uses) receive higher weights than commonly documented ones (e.g., Motivation). Prior empirical work on documentation quality has been limited to manual scoring at small sample sizes (Rondina et al., 2025, n = 100) or single-repository unweighted analyses (Oreamuno et al., 2024, n = 6,758 on HuggingFace Hub only).

Three gaps motivate this work. First, no automated pipeline applies DTS inverse-frequency weighting to documentation scoring. Second, no study operates across multiple repository APIs simultaneously. Third, and less anticipated prior to this work: it is unknown whether the DTS sections designated as highest priority are even representable through structured repository API fields.

We present an automated pipeline that collects structured metadata from HuggingFace Hub (n = 496), OpenML (n = 200), and UCI ML Repository (n = 62), maps API fields to DTS sections, and computes weighted completeness scores. The pipeline was evaluated against two pre-specified gate criteria: corpus coverage rate of at least 0.70 and proxy Pearson r of at least 0.70. Both criteria were met.

The principal empirical finding is a structural asymmetry: the DTS sections with the highest importance weights (Preprocessing at 1.8 and Uses at 1.5) score near-zero across all repositories because these sections exist only as free-text prose in dataset cards, not as structured API fields. Lower-weighted sections such as Motivation (1.0) and Composition (0.9) show substantially higher coverage. This pattern, which we term the documentation API gap, indicates that API-based scoring systematically cannot assess the documentation dimensions that responsible AI frameworks prioritize most.

This paper makes three contributions: (1) the first automated cross-repository DTS-weighted documentation scoring pipeline, achieving 91.8\% coverage across 758 datasets; (2) large-scale empirical characterization of per-section API coverage, revealing the documentation API gap; and (3) identification of the measurement infrastructure needed for future causal studies of documentation quality.

Several important limitations should be noted at the outset. The proxy validation (r = 0.989) measures internal algorithmic consistency between weighted and unweighted scoring, not agreement with human expert judgment. The UCI repository contributed zero scoreable datasets due to a field naming mismatch in the scoring configuration. Two key citations (Rondina et al., 2025; Oreamuno et al., 2024) have not been independently verified in literature databases. These limitations are discussed in detail in Section 6.
